\section*{الفصل \ref{ch:g2:water-on-earth} --- الماء على الأرض}

\begin{solution}{exo:g2:water-on-earth:1}
\omterm{def:g2:water-on-earth:salt-fresh}{ماء مالح}: \omterm{def:g2:water-on-earth:ocean}{البحر}، \omterm{def:g2:water-on-earth:ocean}{والمحيط}. \omterm{def:g2:water-on-earth:salt-fresh}{ماء عذب}: \omterm{def:g1:landscapes-around-us:water}{النهر}، والمطر، \omterm{def:g1:landscapes-around-us:water}{والبحيرة} في \omterm{def:g1:landscapes-around-us:landform}{الجبال}،
والجليد على \omterm{def:g1:landscapes-around-us:landform}{الجبل}.
\end{solution}

\begin{solution}{exo:g2:water-on-earth:2}
لأن معظم سطحها يغطيه ماء المحيطات الأزرق: فإذا نظرنا إليها من الفضاء
بدت زرقاء.
\end{solution}

\begin{solution}{exo:g2:water-on-earth:3}
$100 - 71 = 29$ مربعًا يابسة.
\end{solution}

\begin{solution}{exo:g2:water-on-earth:4}
\omterm{def:g2:water-on-earth:ocean}{المحيط} أكبر: \omterm{def:g2:water-on-earth:ocean}{فالبحر} مساحة أصغر من \omterm{def:g2:water-on-earth:salt-fresh}{الماء المالح}، تقع غالبًا على طرف
\omterm{def:g2:water-on-earth:ocean}{محيط}.
\end{solution}

\begin{solution}{exo:g2:water-on-earth:5}
لا: ماء \omterm{def:g2:water-on-earth:ocean}{المحيط} \omterm{def:g2:water-on-earth:salt-fresh}{ماء مالح}، وشربه يجعل الجسم يفقد الماء. على البحار أن
يحمل معه ماءً عذبًا.
\end{solution}

\begin{solution}{exo:g2:water-on-earth:6}
مثلًا: في \omterm{def:g1:landscapes-around-us:water}{الأنهار}، وفي \omterm{def:g1:landscapes-around-us:water}{البحيرات}، وفي جليد \omterm{def:g1:landscapes-around-us:landform}{الجبال} وثلجها وجليد القارة
القطبية الجنوبية، وتحت الأرض، وفي المطر.
\end{solution}

\begin{solution}{exo:g2:water-on-earth:7}
\omterm{def:g2:water-on-earth:salt-fresh}{الماء المالح}: 96 دلوًا. \omterm{def:g2:water-on-earth:salt-fresh}{الماء العذب}: $2 + 2 = 4$ دلاء.
(وبدقة، جزء من الماء الذي تحت الأرض مالح أيضًا؛ والصورة تعدّه مع \omterm{def:g2:water-on-earth:salt-fresh}{الماء
العذب}.)
\end{solution}

\begin{solution}{exo:g2:water-on-earth:8}
المحيطات (96 دلوًا)، ثم الجليد والثلج (دلوان)، ثم \omterm{def:g1:landscapes-around-us:water}{الأنهار} (أقل بكثير
من دلو واحد).
\end{solution}

\begin{solution}{exo:g2:water-on-earth:9}
قد يتسرب في الأرض؛ وقد يجري على الأرض إلى جدول أو \omterm{def:g1:landscapes-around-us:water}{نهر}؛ وقد يتبخر
عائدًا إلى الهواء حين تظهر الشمس.
\end{solution}

\begin{solution}{exo:g2:water-on-earth:10}
يبقى الملح في قاع الوعاء. والماء تبخر: ذهب إلى الهواء في صورة بخار
ماء.
\end{solution}

\begin{solution}{exo:g2:water-on-earth:11}
مثلًا: تنحدر القطرة على السطح إلى المزراب ثم إلى المصرف، ثم إلى جدول
ثم \omterm{def:g1:landscapes-around-us:water}{نهر}؛ فيحملها \omterm{def:g1:landscapes-around-us:water}{النهر} إلى \omterm{def:g2:water-on-earth:ocean}{البحر}؛ وفي يوم مشمس تتبخر من \omterm{def:g2:water-on-earth:ocean}{البحر}، وتصعد في
الهواء، وتصير جزءًا من غيمة.
\end{solution}
